\chapter{The Integral: Area and Accumulation}

In the previous chapter on Newton's method, the tangent line let us pinpoint
roots of complicated equations using local slopes. For ten chapters, in fact, we
have been obsessed with that one operation: \emph{taking things apart}. The
derivative slices a curve into its instantaneous tendencies -- the slope here,
the slope there, the rate of change at one frozen moment. We zoomed in until the
curve became a line and read off how steep that line was. That is calculus run
\emph{downhill}: from a whole thing to its local behavior.

This chapter runs calculus the other way. We are going to \emph{add things up}.
Given a rate that changes from moment to moment, we will reconstruct the total
that accumulated. Given a curve, we will find the area trapped underneath it.
And the astonishing thing -- the thing the next chapter is entirely about -- is
that these two activities, taking apart and adding up, are exact inverses of one
another. Differentiation undoes integration and integration undoes
differentiation, the way subtraction undoes addition. But we are getting ahead
of ourselves. To appreciate that fundamental connection, you first have to meet the integral on
its own terms, with no derivatives in sight at all.

\begin{focusbox}
Two questions that look unrelated turn out to be the \emph{same} question:
(1) What is the area under a curve? (2) If I know the \emph{rate} at which
something accumulates at every instant, how much accumulates in total? This
chapter builds one object -- the \textbf{definite integral}, measuring the net total accumulated across an interval as a limit of thin slices -- that answers both.
\end{focusbox}

By the end you will know what the symbol $\int_a^b f(x)\,\dd x$ really means: not
as a rule to apply, but as a limit of a perfectly concrete sum of rectangles.
You will see why that elongated S was chosen, what the $\dd x$ is doing there,
and why ``area'' sometimes comes out negative. We will compute an integral or
two from raw first principles -- no shortcuts, no Fundamental Theorem -- just so
you can feel in your hands what the definition is actually asking for.

\chapterline

\section{Two ancient problems, one idea}
\label{ch20:two-problems}

\subsection{The area problem}

How much area is inside a circle? A child can answer for a rectangle (length
times width) or a triangle (half base times height), because those shapes are
built from straight lines and area for straight-line shapes is elementary. But a
circle is bounded by a curve, and there is no finite recipe of multiplications
that gives its area exactly. The number $\pi$ is, at bottom, the confession that
the area of a curved region is a genuinely new kind of quantity.

The Greeks found the way through, and it is the same way we will take. Around
$250$ BCE, Archimedes computed the area of regions bounded by curves -- a
segment of a parabola, the circle -- by trapping the unknown area between shapes
he \emph{could} measure. Inscribe a polygon inside the circle and circumscribe
another outside it. The true area is squeezed between the two. Now add more
sides. The inscribed polygon swells outward, the circumscribed one shrinks
inward, and the gap between them closes. The area of the circle is the single
number both sequences are converging to.

\begin{historybox}
The dream of calculating areas under curves spans over two millennia of struggle against the paradoxes of the infinite. Around 250~BCE, in his treatise \emph{Quadrature of the Parabola} (where ``quadrature'' was the classical term for finding the area of a curved region by matching it to a square), Archimedes calculated the exact area of a parabolic segment. He inscribed a large triangle and packed the remaining gaps with an infinite cascade of smaller triangles whose areas formed a geometric series---a sum where each term is a fixed fraction ($\tfrac{1}{4}$) of the previous one:
\[
  \text{Area} = T + \frac{1}{4}T + \frac{1}{16}T + \cdots = \frac{4}{3}T.
\]
Because ancient Greek mathematics rejected actual infinities, Archimedes avoided limits by using the \emph{method of exhaustion}: a double proof by contradiction showing the area could be neither greater nor less than $\tfrac{4}{3}T$, exhausting any conceivable gap between the polygon and the curve.

Nearly two millennia later, in \emph{Geometria indivisibilibus continuorum} (1635), the Italian mathematician Bonaventura Cavalieri took a daring leap: he imagined a two-dimensional shape as a stack of infinitely many parallel lines of zero thickness, which he called \emph{indivisibles} (like the pages bound in a book, or coins piled in a cylinder). Cavalieri showed that if two regions share equal cross-sections at every height, their total areas must match (\emph{Cavalieri's Principle}). Yet indivisibles sparked a furious controversy. Jesuit mathematicians such as Paul Guldin and the English philosopher Thomas Hobbes attacked them as logically absurd, demanding how lines with zero width could ever accumulate into a finite area, or how an infinity of zeroes could add up to anything at all.

In October 1675, Gottfried Wilhelm Leibniz resolved the conceptual knot with one of the most fruitful inventions in the history of science: the integral sign $\int$. Leibniz realized that instead of zero-width indivisibles, one should sum an infinity of thin rectangles of height $y$ and infinitesimal width $\dd x$---a width smaller than any positive number yet not zero. He adopted the elongated letter ``S'' (for the Latin \emph{summa}, or sum) and wrote $\int y\,\dd x$, so the notation itself explained that an integral is a sum of tiny rectangular strips.

The modern foundation was completed in the nineteenth century by Augustin-Louis Cauchy and Bernhard Riemann. In 1854, Riemann proved that the integral is neither an infinite series of metaphysical indivisibles nor a mysterious sum of zeroes, but the precise limit of finite approximating sums as the mesh size (the width of the widest subinterval) shrinks to zero: $\lim_{\|P\|\to 0} \sum f(x_k^*)\Delta x_k$. This established a universal, rigorous definition of area, volume, and continuous accumulation. Yet Riemann's integral had its own boundary: it fails for wildly discontinuous functions (such as a function that equals $1$ on rational numbers and $0$ on irrationals). In 1902, Henri Lebesgue resolved this by partitioning the function's outputs rather than its inputs, creating modern measure theory. Today, the frontier remains open in mathematical physics: integrating over infinite-dimensional spaces of paths---the Feynman path integrals at the heart of quantum mechanics and quantum field theory---still lacks a universally rigorous mathematical foundation in four spacetime dimensions.
\end{historybox}

\subsection{The accumulation problem}

Now a problem that seems to have nothing to do with area. A car's speedometer
reads $v(t)$ at every instant $t$. The speed is changing -- the driver
accelerates, brakes, coasts. \emph{How far did the car travel} between time
$t=a$ and time $t=b$?

If the speed were \emph{constant} at $60$ miles per hour for $2$ hours, distance
would be easy: $60 \times 2 = 120$ miles. Rate times time. But the speed is not
constant; it changes continuously, and ``rate times time'' has no obvious
meaning when the rate is a moving target.

Here is the method, and notice that it applies the same geometric principle
to motion. Chop the trip into tiny time intervals. Over a
single hundredth of a second the speed barely changes, so on that sliver the
motion is \emph{almost} constant-speed, and distance is \emph{almost} just
``speed at that moment times the length of the sliver.'' Add up the slivers. The
finer you chop, the better the approximation, and the true distance is the
number these sums converge to.

\begin{intuitionbox}
Both problems have the same skeleton. There is a quantity you cannot get
directly because something is \emph{curved} or \emph{varying}. So you cut the
domain into many small pieces, on each piece pretend the varying thing is
constant (which is nearly true when the piece is small), compute the easy
straight-line answer on each piece, add them all up, and then let the pieces
shrink toward zero width. The integral is the limit of that sum.
\end{intuitionbox}

Watch what the distance computation produces. On the time-axis, lay out the
intervals. Above each one, draw a rectangle whose height is the speed during
that interval. The \emph{area} of one rectangle is (height)$\times$(width)
$=$ (speed)$\times$(time) $=$ distance covered in that interval. So the total
distance is the total \emph{area} of all the rectangles -- which is an
approximation to the area under the speed curve $v(t)$. The accumulation problem
\emph{is} the area problem. Distance traveled is the area under the speed graph.

\begin{figure}[h]
\centering
\begin{tikzpicture}
\begin{axis}[
  width=11cm, height=6cm,
  xlabel={time $t$}, ylabel={speed $v(t)$},
  xmin=0, xmax=6, ymin=0, ymax=5,
  xtick={1,2,3,4,5}, ytick=\empty,
  axis lines=left, samples=100, domain=0:6,
  clip=false,
]
\addplot[ybar interval, fill=areafill, fill opacity=0.55, draw=curveblue!60,
  samples at={1,2,3,4,5}] {1.5 + 1.8*sin(deg(0.55*x))};
\addplot[curveblue, very thick, samples=120, domain=0.6:5.4]
  {1.5 + 1.8*sin(deg(0.55*x))};
\node[curveblue] at (axis cs:4.6,4.4) {$v(t)$};
\end{axis}
\end{tikzpicture}
\caption{Distance traveled is the area under the speed curve. Each rectangle has
height equal to the speed on that interval and width equal to the elapsed time,
so its area is a distance. The sum of the rectangles approximates the area under
$v(t)$; the true distance is the limit as the rectangles get thin.}
\label{ch20:fig-speed}
\end{figure}

\section{Riemann sums: making the idea precise}
\label{ch20:riemann}

Let us stop gesturing and pin the construction down. We have a function $f$ on a
closed interval $[a,b]$. For now, picture $f$ as continuous and positive, so
``area under the graph'' is unambiguous -- we will lift those assumptions soon.

\subsection{Partitioning the interval}

Cut $[a,b]$ into $n$ subintervals. Choose points
\[
  a = x_0 < x_1 < x_2 < \cdots < x_{n-1} < x_n = b.
\]
This list is called a \textbf{partition} $P$ of $[a,b]$. In plain words, a partition is simply a collection of slice marks that cut the interval into smaller segments, like knife cuts dividing a loaf of bread. The $k$-th subinterval
is $[x_{k-1}, x_k]$, and its width is
\[
  \Delta x_k = x_k - x_{k-1}.
\]
The widths need not be equal, though most of the time we will make them equal:
if we use $n$ subintervals of the \emph{same} width, then each has width
\[
  \Delta x = \frac{b-a}{n}, \qquad x_k = a + k\,\Delta x.
\]
This even spacing is called a \textbf{regular} (or uniform) partition, and it is
what you should picture by default.

\subsection{Sampling a height on each piece}

On each subinterval $[x_{k-1}, x_k]$ we are going to approximate the curve by a
single horizontal line -- a rectangle. But at what height? The curve takes many
values across the subinterval; we must pick one. Choose any point
$x_k^\ast \in [x_{k-1}, x_k]$, called a \textbf{sample point}, and use the height
$f(x_k^\ast)$. A sample point is our chosen evaluation point inside that slice: we evaluate the function at this single location to fix the height of our approximating rectangle. The rectangle over the $k$-th subinterval then has
\[
  \text{area} = f(x_k^\ast)\,\Delta x_k.
\]
Add them up over all $n$ subintervals.

\begin{definition}[Riemann sum]
\label{ch20:def-riemann-sum}
Given a function $f$ on $[a,b]$, a partition $P$ with subintervals of width
$\Delta x_k$, and a choice of sample points $x_k^\ast \in [x_{k-1},x_k]$, the
corresponding \textbf{Riemann sum} is
\[
  S(P, \{x_k^\ast\}) \;=\; \sum_{k=1}^{n} f(x_k^\ast)\,\Delta x_k.
\]
\end{definition}

This sum is just a number: a finite pile of rectangle areas, signed by the sign
of $f$. It is named for Bernhard Riemann, who in $1854$ gave the first clean
general definition of the integral built exactly this way. Each Riemann sum is
an \emph{approximation} to the area we are after. To get the true area we will
take a limit -- but first let us look at the three most common ways to choose the
sample points, because their pictures teach you what is going on.

\subsection{Left, right, and midpoint sums}
\label{ch20:lrm}

With a regular partition, three natural choices of $x_k^\ast$ stand out.

\begin{itemize}
\item \textbf{Left sum} $L_n$: take $x_k^\ast = x_{k-1}$, the left endpoint of
each subinterval. The rectangle's height is the value of $f$ at the left edge.
\item \textbf{Right sum} $R_n$: take $x_k^\ast = x_k$, the right endpoint.
\item \textbf{Midpoint sum} $M_n$: take $x_k^\ast = \tfrac{1}{2}(x_{k-1}+x_k)$,
the center of each subinterval.
\end{itemize}

\begin{figure}[h]
\centering
\begin{tikzpicture}
\begin{axis}[
title style={font=\small},
  width=5.3cm, height=5cm, title={Left sum $L_4$},
  xmin=0, xmax=2, ymin=0, ymax=4.4, xtick={0,0.5,1,1.5,2},
  ytick=\empty, axis lines=left, domain=0:2, samples=80,
]
\addplot[ybar interval, fill=areafill, fill opacity=0.5, draw=curveblue!60,
  samples at={0,0.5,1,1.5,2}] {x^2+0.3};
\addplot[curveblue, very thick] {x^2+0.3};
\end{axis}
\end{tikzpicture}
\hfill
\begin{tikzpicture}
\begin{axis}[
title style={font=\small},
  width=5.3cm, height=5cm, title={Right sum $R_4$},
  xmin=0, xmax=2, ymin=0, ymax=4.4, xtick={0,0.5,1,1.5,2},
  ytick=\empty, axis lines=left, domain=0:2, samples=80,
]
\addplot[ybar interval, fill=areafill, fill opacity=0.5, draw=curveblue!60,
  samples at={0,0.5,1,1.5,2}] {(x+0.5)^2+0.3};
\addplot[curveblue, very thick] {x^2+0.3};
\end{axis}
\end{tikzpicture}
\hfill
\begin{tikzpicture}
\begin{axis}[
title style={font=\small},
  width=5.3cm, height=5cm, title={Midpoint sum $M_4$},
  xmin=0, xmax=2, ymin=0, ymax=4.4, xtick={0,0.5,1,1.5,2},
  ytick=\empty, axis lines=left, domain=0:2, samples=80,
]
\addplot[ybar interval, fill=areafill, fill opacity=0.5, draw=curveblue!60,
  samples at={0,0.5,1,1.5,2}] {(x+0.25)^2+0.3};
\addplot[curveblue, very thick] {x^2+0.3};
\end{axis}
\end{tikzpicture}
\caption{Three Riemann sums for $f(x)=x^2+0.3$ on $[0,2]$ with four
subintervals. On a rising curve the left sum \emph{underestimates} (each
rectangle ducks below the curve) and the right sum \emph{overestimates}. The
midpoint sum splits the difference and is usually closest.}
\label{ch20:fig-lrm}
\end{figure}

Look at Figure~\ref{ch20:fig-lrm}. On a curve that is \emph{increasing}, the
left endpoint is the lowest point of $f$ on each subinterval, so every left
rectangle sits below the curve and $L_n$ \emph{undershoots} the true area.
Symmetrically, $R_n$ \emph{overshoots}. The true area is sandwiched:
$L_n \le \text{area} \le R_n$. (On a decreasing curve the roles swap.) That
sandwich is Archimedes' inscribed-and-circumscribed idea reborn: the left sum is
an inscribed staircase, the right sum a circumscribed one, and the area lives
between them.

\begin{recoverybox}
\textbf{Pause and think:} What happens if the curve is \emph{decreasing} instead of increasing? Picture it for three seconds before reading on: every left rectangle will now jut above the falling curve, and every right rectangle will duck below it. The inequalities completely flip to $R_n \le \text{area} \le L_n$, but the true area remains securely trapped between the two staircases.
\end{recoverybox}

\begin{labbox}
Let us do one entirely by hand so the definition stops being abstract. Take
$f(x) = x^2$ on $[0,1]$ with $n = 4$ equal subintervals, so $\Delta x =
\tfrac14$ and the partition points are $0, \tfrac14, \tfrac12, \tfrac34, 1$.

\textbf{Left sum.} Heights at $0,\tfrac14,\tfrac12,\tfrac34$:
\[
  L_4 = \tfrac14\!\left(0^2 + \tfrac1{16} + \tfrac14 + \tfrac9{16}\right)
      = \tfrac14 \cdot \tfrac{14}{16} = \tfrac{14}{64} = 0.21875.
\]
\textbf{Right sum.} Heights at $\tfrac14,\tfrac12,\tfrac34,1$:
\[
  R_4 = \tfrac14\!\left(\tfrac1{16} + \tfrac14 + \tfrac9{16} + 1\right)
      = \tfrac14 \cdot \tfrac{30}{16} = \tfrac{30}{64} = 0.46875.
\]
\textbf{Midpoint sum.} Heights at $\tfrac18,\tfrac38,\tfrac58,\tfrac78$:
\[
  M_4 = \tfrac14\!\left(\tfrac1{64} + \tfrac9{64} + \tfrac{25}{64}
        + \tfrac{49}{64}\right) = \tfrac14 \cdot \tfrac{84}{64} = 0.328125.
\]
The true area, as we will prove below, is exactly $\tfrac13 \approx 0.3333$.
Notice $L_4 < \tfrac13 < R_4$ as promised, and $M_4$ is by far the closest --
off by less than $0.006$. Four rectangles already do well.
\end{labbox}

\begin{trapbox}
A Riemann sum is \textbf{not} the integral. It is one finite approximation,
pinned to a particular partition and a particular set of sample points. Change
the sample points and you get a different number. The integral is what all these
numbers \emph{converge to} as the rectangles shrink. Students who write
``$\int_a^b f = \sum f(x_k)\Delta x$'' as if they were equal have skipped the
only step that matters -- the limit. The sum approximates; the limit \emph{is}.
\end{trapbox}

\section{The definite integral as a limit}
\label{ch20:def-integral}

Now we take the limit. As we add more and more subintervals and let their widths
shrink toward zero, the staircase of rectangles hugs the curve ever more tightly,
and all the different Riemann sums -- left, right, midpoint, any sample points at
all -- get squeezed toward one common value. That value is the integral.

We measure the ``fineness'' of a partition by its \textbf{mesh} (or norm), the
width of its widest subinterval:
\[
  \norm{P} = \max_{1 \le k \le n} \Delta x_k.
\]
Think of the mesh as measuring the coarseness of your grid: if even the widest subinterval is squeezed down to paper-thinness, then every single subinterval is forced to vanish as well. Driving $\norm{P} \to 0$ forces \emph{every} subinterval to shrink, rather than merely
the average. With that we can state the definition.

\begin{definition}[The definite integral]
\label{ch20:def-definite-integral}
A function $f$ is \textbf{integrable} on $[a,b]$ if there is a single number $I$
with the following property: the Riemann sums approach $I$ no matter how the
sample points are chosen, as the mesh goes to zero. In everyday terms, being integrable means the curve is well-behaved enough that our rectangle approximations settle down to one unambiguous total area. Precisely, for every
$\eps > 0$ there is a $\delta > 0$ such that for \emph{any} partition $P$ with
$\norm{P} < \delta$ and \emph{any} choice of sample points,
\[
  \abs[\Big]{\sum_{k=1}^{n} f(x_k^\ast)\,\Delta x_k \;-\; I} < \eps.
\]
When this number $I$ exists we call it the \textbf{definite integral} of $f$
from $a$ to $b$ and write
\[
  \int_a^b f(x)\,\dd x \;=\; \lim_{\norm{P}\to 0}
      \sum_{k=1}^{n} f(x_k^\ast)\,\Delta x_k \;=\; I.
\]
\end{definition}

\begin{figure}[H]
\centering
\begin{tikzpicture}
% Panel 1: n = 4
\begin{axis}[
  name=p1,
  width=5.1cm, height=4.6cm,
  title style={font=\small},
  title={$n=4$ subintervals},
  xmin=0, xmax=2.0, ymin=0, ymax=4.2,
  xtick={0, 1, 2}, ytick=\empty,
  axis lines=left,
]
\addplot[ybar interval, fill=areafill, fill opacity=0.45, draw=curveblue!70,
  samples at={0, 0.5, 1.0, 1.5, 2.0}] {0.3 + 0.85*x^2};
\addplot[curveblue, very thick, domain=0:2, samples=60] {0.3 + 0.85*x^2};
\node[font=\footnotesize] at (axis cs:1.25, 0.5) {$S_4$};
\end{axis}

% Panel 2: n = 16
\begin{axis}[
  at={(p1.east)}, anchor=west, xshift=0.35cm,
  name=p2,
  width=5.1cm, height=4.6cm,
  title style={font=\small},
  title={$n=16$ subintervals},
  xmin=0, xmax=2.0, ymin=0, ymax=4.2,
  xtick={0, 1, 2}, ytick=\empty,
  axis lines=left,
]
\addplot[ybar interval, fill=areafill, fill opacity=0.45, draw=curveblue!70,
  samples at={0, 0.125, 0.25, 0.375, 0.5, 0.625, 0.75, 0.875, 1.0, 1.125, 1.25, 1.375, 1.5, 1.625, 1.75, 1.875, 2.0}] {0.3 + 0.85*x^2};
\addplot[curveblue, very thick, domain=0:2, samples=60] {0.3 + 0.85*x^2};
\node[font=\footnotesize, fill=white, inner sep=1.5pt, rounded corners=1pt] at (axis cs:1.0, 0.5) {$S_{16}$};
\end{axis}

% Panel 3: n -> infinity with differential strip
\begin{axis}[
  at={(p2.east)}, anchor=west, xshift=0.35cm,
  name=p3,
  width=5.1cm, height=4.6cm,
  title style={font=\small},
  title={$n\to\infty$: $\int_a^b f(x)\,\dd x$},
  xmin=0, xmax=2.0, ymin=0, ymax=4.2,
  xtick={0, 1, 2}, ytick=\empty,
  axis lines=left,
  clip=false,
]
\addplot[fill=areafill, fill opacity=0.35, draw=none, domain=0:2, samples=60] {0.3 + 0.85*x^2} \closedcycle;
\addplot[curveblue, very thick, domain=0:2, samples=60] {0.3 + 0.85*x^2};
\fill[structurepurple!35, draw=structurepurple, thick] (axis cs:1.1,0) rectangle (axis cs:1.25, 1.45);
\node[structurepurple!80!black, font=\tiny, anchor=east] at (axis cs:0.95, 2.3) {$\dd A = f(x)\,\dd x$};
\draw[->, >=stealth, structurepurple] (axis cs:1.0, 2.2) -- (axis cs:1.175, 1.48);
\draw[<->, >=stealth, structurepurple] (axis cs:1.1, -0.25) -- (axis cs:1.25, -0.25)
  node[midway, below, font=\tiny] {$\dd x$};
\end{axis}
\end{tikzpicture}
\caption{The limit of Riemann sums as the partition refines.
Left: A coarse partition ($n=4$) leaves visible staircase gaps and overshoots.
Center: A finer partition ($n=16$) hugs the curve closely as the staircase smooths out.
Right: In the continuous limit $n\to\infty$ ($\Delta x \to 0$), the sum becomes the
definite integral $\int_a^b f(x)\,\dd x$, accumulating infinitely many infinitesimal strips
$\dd A = f(x)\,\dd x$.}
\label{ch20:fig-limit-progression}
\end{figure}

That ``for any choice of sample points'' is the whole strength of the
definition. It is demanding that left sums, right sums, midpoint sums, and every
other choice all funnel to the \emph{same} limit. For a badly
behaved $f$ they might not, and then $f$ is not integrable. The foundational fact
of the subject guarantees that most familiar functions qualify:

\begin{theorem}[Integrability of continuous functions]
\label{ch20:thm-cont-integrable}
If $f$ is continuous on $[a,b]$, then $f$ is integrable on $[a,b]$. More
generally, $f$ is integrable if it is bounded and has at most finitely many
discontinuities.
\end{theorem}

The reason continuity is enough is intuitive: on a tiny subinterval a continuous
function cannot wiggle much (that is essentially \emph{uniform} continuity on the
closed interval $[a,b]$---the mathematical guarantee that on a closed interval, the curve's variation across any slice shrinks uniformly to zero as the slice narrows), so the gap between the tallest and shortest rectangle
on that piece is small, which means left and right sums cannot disagree by much,
which means everything is trapped together. Squeeze the trap shut and the limit
exists.

\subsection{Decoding the notation}

The symbol $\int_a^b f(x)\,\dd x$ is one of the best-designed notations in all of
mathematics, and Leibniz designed it in $1675$ to be read literally.

\begin{happeningbox}
Stare at the symbol and read what each part was \emph{built} to mean.
\begin{itemize}
\item In the expression $f(x)\,\dd x$, the factor $\dd x$ recalls the finite subinterval width $\Delta x_k$ as we pass to the limit $\Delta x_k \to 0$.
\item $f(x)$ is the height of the curve across that slice.
\item The product $f(x)\,\dd x$ represents the area of one approximating strip: height times width.
\item The sign $\int$ is an elongated letter $S$, for \emph{summa} -- ``sum.''
Leibniz stretched the symbol $\Sigma$ into a continuous $S$
because we are no longer adding finitely many terms, but passing to a continuous limit.
\item The $\int_a^b$ says: accumulate these contributions as $x$ sweeps from $a$ to $b$.
\end{itemize}
Read it aloud the way Leibniz meant it: \emph{``sum, from $a$ to $b$, of height
times infinitesimal width.''} The integral is, in one symbol, the instruction
``add up infinitely many infinitely thin pieces.'' Everything we did with
Riemann sums was the finite, limit-respecting version of exactly that
sentence.
\end{happeningbox}

A word on the $x$. It is a \textbf{dummy variable}: it ranges over the
subintervals and then disappears in the limit. A dummy variable is simply a temporary placeholder whose label never affects the final outcome, exactly like the index $k$ in a summation $\sum_{k=1}^n k$. The value of $\int_a^b f(x)\,\dd
x$ is a fixed number that does not depend on what letter you used. So
$\int_a^b f(x)\,\dd x = \int_a^b f(t)\,\dd t = \int_a^b f(u)\,\dd u$. The
variable is scaffolding; the number is the building.

\begin{whybox}
Why does letting the rectangle width go to zero give the \emph{exact} area
rather than just a better guess? Because of the squeeze. On any subinterval, the
true area under the curve lies between the shortest rectangle (height $=$ minimum
of $f$ there) and the tallest (height $=$ maximum). Summed up, the true area lies
between the lower sum $\underline{S}_n$ and the upper sum $\overline{S}_n$. Each
individual gap is at most (oscillation of $f$ on that piece)$\times$(width of
that piece). As the mesh shrinks, continuity forces every oscillation toward
zero, so the total gap $\overline{S}_n - \underline{S}_n$ collapses to zero. The
true area is pinned between two quantities that are coming together: it can be
nothing but their shared limit. No error survives the limit -- that is the whole
trick of calculus, and it is Archimedes' method of exhaustion made rigorous.
\end{whybox}

\section{Computing an integral from the definition}
\label{ch20:from-def}

Before we have any shortcuts, let us earn one integral the hard way -- straight
from Definition~\ref{ch20:def-definite-integral} -- so you trust that the limit
really is there and really is computable. We will compute
\[
  \int_0^1 x^2 \, \dd x
\]
and find it equals $\tfrac13$, the number we sampled in the lab box above.

\begin{example}[$\int_0^1 x^2\,\dd x$ by brute force]
\label{ch20:ex-x2}
Use a regular partition into $n$ equal pieces, so $\Delta x = \tfrac1n$ and
$x_k = \tfrac{k}{n}$. Choose right endpoints as sample points, $x_k^\ast =
\tfrac{k}{n}$. The Riemann sum is
\[
  R_n = \sum_{k=1}^{n} \left(\frac{k}{n}\right)^2 \cdot \frac{1}{n}
      = \frac{1}{n^3} \sum_{k=1}^{n} k^2.
\]
Now we need the sum of the first $n$ squares -- a classic closed form,
provable by induction:
\[
  \sum_{k=1}^{n} k^2 = \frac{n(n+1)(2n+1)}{6}.
\]
Substitute it in:
\[
  R_n = \frac{1}{n^3} \cdot \frac{n(n+1)(2n+1)}{6}
      = \frac{(n+1)(2n+1)}{6n^2}
      = \frac{2n^2 + 3n + 1}{6n^2}.
\]
Divide through by $n^2$:
\[
  R_n = \frac{1}{6}\left(2 + \frac{3}{n} + \frac{1}{n^2}\right).
\]
Now let $n \to \infty$ (which drives the mesh $\tfrac1n \to 0$). The
$\tfrac{3}{n}$ and $\tfrac{1}{n^2}$ terms vanish, leaving
\[
  \int_0^1 x^2 \,\dd x = \lim_{n\to\infty} R_n
      = \frac{1}{6}\cdot 2 = \frac{1}{3}.
\]
\end{example}

There it is, exactly $\tfrac13$, with no Fundamental Theorem, no antiderivative,
nothing but a sum of squares and a limit. Notice the structure of the
computation: a finite-but-growing sum, a closed form for it, then a limit that
washes away every term except the one that survives. That is the integral doing
precisely what its definition promised.

\begin{labbox}
Repeat the computation with \emph{left} endpoints, $x_k^\ast = \tfrac{k-1}{n}$.
Using $\sum_{k=1}^n (k-1)^2 = \sum_{j=0}^{n-1} j^2 = \frac{(n-1)n(2n-1)}{6} = \frac{2n^3 - 3n^2 + n}{6}$, you will find $L_n = \frac{1}{6}\left(2 - \frac{3}{n} + \frac{1}{n^2}\right)$. This value is \emph{below} $\tfrac13$ for every finite $n$ (the left sum undershoots a
rising curve, as Figure~\ref{ch20:fig-lrm} showed) -- yet it still limits to the
same $\tfrac13$. Left and right disagree at every finite stage and agree in the
limit. That agreement is integrability in action.
\end{labbox}

\begin{aestheticbox}
Every finite sum $R_n$ was slightly too big because the rectangles poked above the parabola, yet as $n\to\infty$ the excess $\tfrac{3}{6n} + \tfrac{1}{6n^2}$ melted away completely, leaving the exact area $\tfrac13$. Finite approximations plus a limit gave an exact curved area using nothing more than algebra.
\end{aestheticbox}

\section{Signed area: when the curve dips below the axis}
\label{ch20:signed-area}

So far $f$ was positive and ``area under the curve'' meant exactly that. But the
definition of the integral never required $f \ge 0$. If $f(x_k^\ast)$ is
\emph{negative}, the term $f(x_k^\ast)\,\Delta x_k$ is negative -- it
\emph{subtracts} from the sum. So what does the integral compute when the curve
goes below the $x$-axis?

\begin{focusbox}
The definite integral computes \textbf{signed area}: area above the $x$-axis
counts as positive, area below counts as negative. The integral is the area
above minus the area below.
\end{focusbox}

\begin{figure}[h]
\centering
\begin{tikzpicture}
\begin{axis}[
  width=11cm, height=5.5cm,
  xlabel={$x$}, ylabel={$f(x)$},
  xmin=-0.3, xmax=6.4, ymin=-2.6, ymax=2.6,
  xtick={1,2,3,4,5}, ytick={-2,-1,1,2},
  axis lines=middle, samples=120, domain=0:6,
]
\addplot[name path=f, curveblue, very thick] {1.8*sin(deg(x))};
\path[name path=axis] (axis cs:0,0) -- (axis cs:6,0);
\addplot[targetgreen!55, fill opacity=0.5] fill between[of=f and axis,
  soft clip={domain=0:3.1416}];
\addplot[trapred!45, fill opacity=0.5] fill between[of=f and axis,
  soft clip={domain=3.1416:6}];
\draw[gray, dashed] (axis cs:6,0) -- (axis cs:6,-0.5);
\node[above, font=\footnotesize] at (axis cs:6,0.06) {$6$};
\node[targetgreen!60!black] at (axis cs:1.55,0.85) {$+$};
\node[trapred!75!black] at (axis cs:4.7,-0.85) {$-$};
\end{axis}
\end{tikzpicture}
\caption{Signed area. For $f(x)=1.8\sin x$ on $[0,6]$, the green hump above the
axis counts positively and the red dip below counts negatively. The integral
$\int_0^6 1.8\sin x\,\dd x$ is the green area minus the red area, not the total
amount of paint.}
\label{ch20:fig-signed}
\end{figure}

This sign convention is not an arbitrary bookkeeping rule; it is forced on us,
and it is exactly what we \emph{want}. Recall the car. If $v(t)$ is the car's
\emph{velocity} (a signed quantity -- positive going forward, negative in
reverse), then $\int_a^b v(t)\,\dd t$ is the \emph{displacement}: net change in
position. Time spent in reverse should subtract from how far you have gotten, and
the negative integrand makes it do exactly that. If instead you want total
distance traveled -- odometer mileage, which never decreases -- you integrate the
\emph{speed} $\abs{v(t)}$, which is never negative. Displacement and distance
differ precisely by the sign convention, and both are useful.

\begin{recoverybox}
\textbf{Pause and think:} If a car drives east at $50\text{ mph}$ for an hour, then reverses west at $30\text{ mph}$ for an hour, what is its net displacement? The signed integral gives $50(1) + (-30)(1) = 20\text{ miles}$ east. But what distance did the tires roll? The odometer integrates speed $\abs{v(t)}$: $50(1) + 30(1) = 80\text{ miles}$. The definite integral tracks where you ended up; integrating the absolute value tracks the total tire wear.
\end{recoverybox}

\begin{trapbox}
``Find the area between the curve and the axis'' and ``evaluate the integral''
are different requests when the curve crosses the axis. The integral gives signed
area and can be zero even when there is plenty of region (a symmetric wiggle
cancels itself: $\int_0^{2\pi}\sin x\,\dd x = 0$). For \emph{geometric} area --
total paint -- you must split at the crossing points and integrate $\abs{f}$, or
equivalently add up the absolute values of the pieces. Read the question
carefully.
\end{trapbox}

\section{Properties of the integral}
\label{ch20:properties}

Working from the definition every time would be unbearable. Fortunately the
integral inherits a handful of clean properties straight from the Riemann sum --
each one is just a fact about sums that survives the limit. Let $f$ and $g$ be
integrable on the relevant intervals and $c$ a constant.

\begin{focusbox}
\begin{enumerate}[label=\textbf{(\arabic*)}]
\item \textbf{Linearity (sum).}\quad
  $\displaystyle \int_a^b \big(f(x)+g(x)\big)\,\dd x
   = \int_a^b f(x)\,\dd x + \int_a^b g(x)\,\dd x.$
\item \textbf{Linearity (constant multiple).}\quad
  $\displaystyle \int_a^b c\,f(x)\,\dd x = c\int_a^b f(x)\,\dd x.$
\item \textbf{Additivity over intervals.}\quad For $a < c < b$,
  $\displaystyle \int_a^b f = \int_a^c f + \int_c^b f.$
\item \textbf{Orientation (reversing the limits).}\quad
  $\displaystyle \int_b^a f = -\int_a^b f,$ \quad and \quad
  $\displaystyle \int_a^a f = 0.$
\item \textbf{Monotonicity.}\quad If $f(x) \le g(x)$ on $[a,b]$, then
  $\displaystyle \int_a^b f \le \int_a^b g.$
\item \textbf{The constant.}\quad
  $\displaystyle \int_a^b c \,\dd x = c\,(b-a).$
\end{enumerate}
\end{focusbox}

Each of these is obvious once you see it as a statement about the rectangles.

\textbf{Why linearity holds.} A Riemann sum for $f+g$ is
$\sum \big(f(x_k^\ast)+g(x_k^\ast)\big)\Delta x_k$. Finite sums split across
addition, so this is $\sum f(x_k^\ast)\Delta x_k + \sum g(x_k^\ast)\Delta x_k$ --
a Riemann sum for $f$ plus a Riemann sum for $g$, on the \emph{same} partition
and \emph{same} sample points. Take the limit of both sides; the limit of a sum
is the sum of the limits. Pulling a constant $c$ out is the same story: $\sum
c\,f(x_k^\ast)\Delta x_k = c\sum f(x_k^\ast)\Delta x_k$. Linearity is nothing but
the distributive law surviving a limit.

\textbf{Why additivity holds.} If the splitting point $c$ is one of your
partition points, then the rectangles over $[a,b]$ are literally the rectangles
over $[a,c]$ followed by the rectangles over $[c,b]$. The sum over the whole is
the sum over the left part plus the sum over the right part. Take the limit.
Figure~\ref{ch20:fig-additivity} illustrates the geometry: dividing the interval
at $c$ cuts the total region into two adjacent slices that fit together edge to
edge. This is the property that lets us handle a function defined piecewise, or
split at a sign change to compute geometric area.

\begin{figure}[h]
\centering
\begin{tikzpicture}
\begin{axis}[
  width=10.5cm, height=4.8cm,
  xlabel={$x$}, ylabel={$f(x)$},
  xmin=0.5, xmax=5.5, ymin=0, ymax=3.5,
  xtick={1, 3, 5}, xticklabels={$a$, $c$, $b$}, ytick=\empty,
  axis lines=left, samples=100, domain=0.8:5.2,
]
\addplot[name path=f, curveblue, very thick] {1.4 + 0.9*sin(deg(0.9*x - 0.5))};
\path[name path=base] (axis cs:1,0) -- (axis cs:5,0);
\addplot[areafill, fill opacity=0.4] fill between[of=f and base, soft clip={domain=1:3}];
\addplot[targetgreen!60, fill opacity=0.4] fill between[of=f and base, soft clip={domain=3:5}];
\draw[dashed, gray!80!black, thick] (axis cs:3,0) -- (axis cs:3, 2.13);
\node[sourceblue!80!black, font=\small] at (axis cs:2.0, 0.8) {Area $A_1$};
\node[targetgreen!80!black, font=\small] at (axis cs:4.0, 0.8) {Area $A_2$};
\node[curveblue, font=\small, anchor=south] at (axis cs:2.2, 2.35) {$f(x)$};
\end{axis}
\end{tikzpicture}
\caption{Additivity over intervals. Slicing the domain $[a,b]$ at $x=c$ cuts the region into two adjacent pieces of areas $A_1 = \int_a^c f(x)\,\dd x$ and $A_2 = \int_c^b f(x)\,\dd x$. The total area is the sum of the parts: $\int_a^b f(x)\,\dd x = A_1 + A_2 = \int_a^c f(x)\,\dd x + \int_c^b f(x)\,\dd x$.}
\label{ch20:fig-additivity}
\end{figure}

\textbf{Why orientation flips the sign.} We built the integral with $a$ on the
left and $b$ on the right, so $\Delta x_k > 0$. If we agree to sweep the other
way, every $\Delta x_k$ becomes negative, flipping the sign of the whole sum.
Defining $\int_b^a f = -\int_a^b f$ makes the additivity rule
$\int_a^c + \int_c^b = \int_a^b$ hold for \emph{any} arrangement of $a,b,c$ --
even when $c$ is outside $[a,b]$. The convention is chosen to keep the algebra
consistent without splitting into cases, and you will be grateful for it in the
next chapter.

\begin{intuitionbox}
Monotonicity is the most intuitive of all: if one curve sits above another
everywhere, the area under it is larger. There is simply more region underneath.
In rectangle terms, $f \le g$ makes every rectangle for $f$ no taller than the
corresponding rectangle for $g$, so the sums obey the same inequality, and the
inequality passes to the limit. From this one fact flow all the comparison and
bounding estimates you will ever make for integrals.
\end{intuitionbox}

\begin{example}[Properties in action]
\label{ch20:ex-properties}
Suppose someone tells you $\int_0^1 x^2\,\dd x = \tfrac13$ (we proved it) and
$\int_0^1 x\,\dd x = \tfrac12$ (the area of a triangle of base and height $1$).
Then with no further integration:
\[
  \int_0^1 \big(6x^2 - 4x + 5\big)\,\dd x
  = 6\!\int_0^1\! x^2\,\dd x - 4\!\int_0^1\! x\,\dd x + \int_0^1\! 5\,\dd x
  = 6\cdot\tfrac13 - 4\cdot\tfrac12 + 5\cdot 1 = 2 - 2 + 5 = 5.
\]
Linearity and the constant rule did all the work, dissolving a complicated
integrand into simple known pieces.
\end{example}

\section{The average value of a function}
\label{ch20:average}

Here is a question the integral handles cleanly, and one that previews ideas to
come. The average of finitely many numbers is their sum divided by how
many there are. But what is the average \emph{height} of a curve over an interval
-- the average of \emph{infinitely many} values $f(x)$ as $x$ ranges over
$[a,b]$? You cannot divide ``the sum of infinitely many values'' by ``infinity.''
The integral steps in exactly where the naive average breaks down.

Sample the curve at $n$ evenly spaced points and average those:
\[
  \text{average} \approx \frac{f(x_1)+f(x_2)+\cdots+f(x_n)}{n}.
\]
Multiply and divide by $\Delta x = \tfrac{b-a}{n}$. Since $\tfrac1n =
\tfrac{\Delta x}{b-a}$,
\[
  \frac{1}{n}\sum_{k=1}^{n} f(x_k)
  = \frac{1}{b-a}\sum_{k=1}^{n} f(x_k)\,\Delta x.
\]
The sum on the right is a Riemann sum for $f$. Let $n \to \infty$ and it becomes
an integral, giving us the definition.

\begin{definition}[Average value]
\label{ch20:def-average}
The \textbf{average value} of an integrable function $f$ over $[a,b]$ is
\[
  \bar f = \frac{1}{b-a}\int_a^b f(x)\,\dd x.
\]
\end{definition}

\begin{intuitionbox}
Rearrange it as $\bar f \cdot (b-a) = \int_a^b f(x)\,\dd x$. The left side is the
area of a \emph{rectangle} of width $(b-a)$ and height $\bar f$. So the average
value is the height of the rectangle that has the \emph{same area} as the region
under the curve -- the level you would pour the area into to make it flat. If you
flooded the area under the graph and let it settle, $\bar f$ is the water line.
\end{intuitionbox}

\begin{figure}[h]
\centering
\begin{tikzpicture}
\begin{axis}[
  width=10cm, height=5cm,
  xlabel={$x$}, ymin=0, ymax=3.4, xmin=0, xmax=3,
  xtick={0,3}, xticklabels={$a$,$b$}, ytick={1.5}, yticklabels={$\bar f$},
  axis lines=left, samples=120, domain=0:3,
]
\addplot[name path=f, curveblue, very thick] {1.5 + 1.1*sin(deg(1.6*x-1))};
\path[name path=base] (axis cs:0,0) -- (axis cs:3,0);
\addplot[areafill, fill opacity=0.35] fill between[of=f and base];
\addplot[tangentred, very thick, dashed, domain=0:3] {1.805};
\node[tangentred, font=\small, anchor=south] at (axis cs:2.45, 1.88) {$y=\bar f$};
\end{axis}
\end{tikzpicture}
\caption{Average value as a leveling line. The dashed line $y=\bar f$ is set at
the height that makes the rectangle over $[a,b]$ have exactly the same area as
the shaded region under $f$. The curve's excess above the line exactly fills its
deficits below.}
\label{ch20:fig-average}
\end{figure}

\begin{example}[Average value of a parabola]
\label{ch20:ex-avg}
The average value of $f(x)=x^2$ on $[0,1]$ is
\[
  \bar f = \frac{1}{1-0}\int_0^1 x^2\,\dd x = \frac{1}{1}\cdot\frac13 = \frac13.
\]
Sanity check: $f$ runs from $0$ up to $1$ on this interval, but it spends most of
the interval near the bottom (it only climbs steeply near $x=1$). Its values near
the right end, all the way up to $1$, pull the average above the midpoint value
$f(\tfrac12)=\tfrac14$, landing at $\tfrac13$ -- comfortably between the
function's low values and its high value of $1$, just as an average should be.
\end{example}

\begin{connectionbox}
The average value is the bridge to the \emph{Mean Value Theorem for integrals},
which you will meet next chapter: if $f$ is continuous, it actually \emph{attains}
its average value at some point $c\in[a,b]$, i.e. $f(c)=\bar f$. Geometrically,
the curve must cross its own leveling line somewhere. This is the integral's echo
of the Mean Value Theorem you proved for derivatives, and it is one of the gears
that drives the Fundamental Theorem of Calculus.
\end{connectionbox}

\begin{applicationbox}
Average value is everywhere once you look. The root-mean-square voltage of
household AC power is built from the average value of $V(t)^2$ over a cycle -- it
is why a ``$120$-volt'' outlet has a peak voltage near $170$. In probability, the
\emph{expected value} of a continuous random variable is $\int x\,p(x)\,\dd x$, an
average of outcomes weighted by the density $p$. In thermodynamics, the average
energy of a system is an integral over all its states. Whenever you must
summarize a continuously varying quantity by a single representative number, the
average value -- an integral divided by a length -- is the right way to do it.
\end{applicationbox}

\section{Accumulation: the integral as a running total}
\label{ch20:accumulation}

We close by returning to the theme that names this chapter, because it is the
viewpoint that makes the next chapter inevitable. We have treated $\int_a^b f$ as
a fixed number -- one area. But let the \emph{upper limit} move. Define a new
function
\[
  A(x) = \int_a^x f(t)\,\dd t,
\]
the accumulated area from the fixed start $a$ out to a variable endpoint $x$.
(Note the dummy variable is now $t$, freeing $x$ to be the endpoint.) As $x$
slides right, $A(x)$ grows by sweeping in more area; this $A$ is the
\textbf{accumulation function}. Think of it as a running meter or odometer: given any chosen endpoint $x$, it reports the grand total accumulated from the starting point $a$ up to $x$.

\begin{happeningbox}
Think of $f$ as a \emph{rate} and $A$ as the \emph{running total} that rate
builds up. If $f(t)$ is the rate at which water flows into a tank (liters per
second), then $A(x)=\int_a^x f(t)\,\dd t$ is the total liters that have flowed in
by time $x$. If $f(t)$ is a country's birth rate (people per year), $A$ is the
cumulative number of births. If $f$ is velocity, $A$ is displacement. The
integrand is always a \emph{rate of accumulation}; the integral is always the
\emph{amount accumulated}. This single picture -- rate in, total out -- is the
beating heart of integral calculus, and it is the exact inverse of the
derivative's picture (total in, rate out).
\end{happeningbox}

\begin{figure}[h]
\centering
\begin{tikzpicture}
\begin{axis}[
  name=top,
  width=11cm, height=4cm,
  ylabel={rate $f(t)$},
  xmin=0, xmax=4, ymin=0, ymax=3,
  xtick={1,2,3,4}, ytick=\empty,
  axis lines=left, samples=100, domain=0:4, clip=false,
]
\addplot[name path=f, curveblue, very thick] {1.2 + 1.1*sin(deg(1.3*x))};
\path[name path=base] (axis cs:0,0) -- (axis cs:4,0);
\addplot[areafill, fill opacity=0.35] fill between[of=f and base,
  soft clip={domain=0:2.6}];
\draw[dashed, gray!80!black, thick] (axis cs:2.6,0) -- (axis cs:2.6, 0.94);
\fill[curveblue] (axis cs:2.6, 0.94) circle (2pt);
\node[curveblue, font=\footnotesize, anchor=south west] at (axis cs:2.6, 0.94) {$f(x)$};
\node[curveblue] at (axis cs:3.6,2.55) {$f(t)$};
\node[anchor=north] at (axis cs:2.6,-0.05) {$x$};
\end{axis}
\begin{axis}[
  at={(top.below south west)}, anchor=north west, yshift=-0.5cm,
  width=11cm, height=4cm,
  xlabel={endpoint $x$}, ylabel={total $A(x)$},
  xmin=0, xmax=4, ymin=0, ymax=6,
  xtick={1,2,3,4}, ytick=\empty,
  axis lines=left, samples=100, domain=0:4, clip=false,
]
\addplot[targetgreen!70!black, very thick]
  {1.2*x + 0.8462*(1 - cos(deg(1.3*x)))};
\draw[dashed, gray!80!black, thick] (axis cs:2.6,0) -- (axis cs:2.6, 4.78);
\fill[targetgreen!70!black] (axis cs:2.6, 4.78) circle (2pt);
\node[targetgreen!70!black, font=\footnotesize, anchor=south east] at (axis cs:2.45, 4.85) {$A(x)$};
\node[anchor=north] at (axis cs:2.6,-0.05) {$x$};
\end{axis}
\end{tikzpicture}
\caption{The accumulation function $A(x)=\int_a^x f(t)\,\dd t$ (bottom) records
the running area swept out under the rate $f$ (top) as the endpoint $x$ moves
right. Where $f$ is tall, $A$ climbs steeply; where $f$ is small, $A$ nearly
levels off. The \emph{slope} of the lower curve at each $x$ is precisely the
\emph{height} of the upper curve at that same $x$ -- a picture of
$A'(x)=f(x)$ before we prove it.}
\label{ch20:fig-accumulation}
\end{figure}

Now ask the question that ties the whole subject together: \emph{how fast does the
running total $A(x)$ grow?} When the endpoint nudges forward by a tiny $\dd x$,
the area gains a thin sliver of width $\dd x$ and height $f(x)$ -- so
$A$ grows by about $f(x)\,\dd x$, which means
\[
  \dv{A}{x} \;\approx\; f(x).
\]
The \emph{rate of accumulation of the total is the integrand itself}. The
derivative of the area function is the curve. You can see it in
Figure~\ref{ch20:fig-accumulation}: the running-total curve rises steeply
where the rate is large and flattens where the rate is small, so its
slope matches the height of $f$ point for point. Differentiation
and integration are inverse operations. That relationship, made exact, is the
Fundamental Theorem of Calculus -- the subject of the next chapter. We
built the integral from sums and limits alone;
and yet at the end the derivative reappears as the rate of change of accumulated area.

\begin{walkbox}
Take this on your walk. You have known since childhood that distance is speed
times time -- when the speed is constant. The integral is the precise meaning of
``speed times time'' when the speed is \emph{not} constant: chop time fine, on
each sliver the speed is nearly constant so distance is nearly speed-times-sliver,
add up the slivers, take the limit. Now turn it around: if you only had a record
of total distance versus time, how would you recover the speed at an instant? You
would differentiate. So the speedometer and the odometer are derivative and
integral of one another. Carry that one example with you; the Fundamental Theorem
is just this sentence, made airtight, for \emph{every} pair of rate and total.
\end{walkbox}

\section{Looking back, and looking ahead}
\label{ch20:summary}

We built the definite integral from the ground up. We took two ancient,
seemingly separate problems -- the area under a curve and the total accumulated
from a varying rate -- and discovered they are the same problem. We solved it
Archimedes' way: cut the domain into thin pieces, approximate the varying thing
by a constant on each piece, sum the easy answers, and take a limit as the pieces
vanish. That limit, when it exists, is the Riemann integral
$\int_a^b f(x)\,\dd x$, and we proved by hand that $\int_0^1 x^2\,\dd x =
\tfrac13$ with nothing but a sum of squares.

We saw that the integral computes \emph{signed} area, that it obeys clean
linearity, additivity, orientation, and monotonicity rules -- each one a fact
about finite sums that survives the limit -- and that dividing it by the interval
length gives the average value of a function. Finally, by letting the upper limit
roam, we built the accumulation function and glimpsed the equation
$\dv{A}{x}=f(x)$ that will, next chapter, unify the two halves of calculus into a
single theorem.

\begin{aestheticbox}
The derivative tore a curve down to its instantaneous slopes. The integral
rebuilds a total from its instantaneous rates. One is analysis, the other is
synthesis; one takes apart, the other puts together. That they are exact inverses
is not a convenience but a deep structural truth about change itself -- and it is
the reason a single subject, calculus, can describe both how things vary and how
much accumulates. You have now met both halves separately. Turn the page and
watch them fuse.
\end{aestheticbox}

\subsection{The Scope of Accumulation and the Bridge Ahead}
\label{ch20:what-unlocked}

\begin{connectionbox}
The definite integral provides a unified foundation for physical accumulation:
\begin{itemize}
  \item \emph{Continuous accumulation}: Distance traveled under variable velocity ($s = \int v(t)\,\dd t$), total work done by a variable force ($W = \int F(x)\,\dd x$), and mass under non-uniform density are all unified under a single operation.
  \item \emph{Signed net change}: By respecting the algebraic sign of $f(x)$, the integral captures net balances, cancellation between positive inflows and negative outflows, and physical displacements.
  \item \emph{Rigorous measure of curved geometry}: Regions bounded by non-linear functions have unambiguous, mathematically well-defined areas as the limit of Riemann sums.
\end{itemize}

Yet Riemann sums are better as a \emph{definition} than as a daily \emph{calculation tool}. To compute $\int_a^b f(x)\,\dd x$ directly from the definition:
\[
  \int_a^b f(x)\,\dd x = \lim_{n\to\infty} \sum_{k=1}^n f(x_k^*)\,\Delta x,
\]
one needs closed-form formulas for the finite sums $\sum_{k=1}^n f(x_k^*)$. We know such formulas for low-degree polynomials ($\sum k, \sum k^2, \sum k^3$), but for almost any other function---$\sin x$, $e^x$, $\sqrt{1-x^2}$, $\ln x$---there are no simple summation formulas.

The practical shortcut arrives in Chapter~21 (\emph{The Fundamental Theorem of Calculus}). Isaac Newton and Gottfried Wilhelm Leibniz recognized that accumulation and differentiation undo each other:
\[
  \frac{\dd}{\dd x} \int_a^x f(t)\,\dd t = f(x).
\]
Instead of constructing rectangles and taking infinite limits of sums, computing $\int_a^b f(x)\,\dd x$ requires only finding an \emph{antiderivative} $F$ with $F'(x) = f(x)$, and evaluating $F(b) - F(a)$.
\end{connectionbox}

\section{Exercises}
\label{ch20:exercises}

\begin{exercise}[Mental checks: signs, symmetry, and constants]
Answer each of the following in your head in ten seconds, using geometry and symmetry rather than algebraic calculation:
\begin{enumerate}[label=\textbf{(\alph*)}]
\item What is $\int_{-3}^3 x^5\,\dd x$, and why?
\item What is $\int_2^7 4\,\dd x$? (Picture the rectangle.)
\item If $\int_0^5 f(x)\,\dd x = 12$ and $\int_0^3 f(x)\,\dd x = 8$, what is $\int_3^5 f(x)\,\dd x$?
\item Without computing, explain why $\int_0^{\pi} \sin x\,\dd x > 0$ while $\int_0^{2\pi} \sin x\,\dd x = 0$.
\end{enumerate}
\end{exercise}

\begin{exercise}[Hand-computed Riemann sums]
For $f(x) = \sqrt{x}$ on $[0,4]$ with $n=4$ equal subintervals
($\Delta x = 1$), compute the left sum $L_4$, the right sum $R_4$, and the
midpoint sum $M_4$. The exact value of $\int_0^4 \sqrt{x}\,\dd x$ is
$\tfrac{16}{3}\approx 5.333$. Which of your three sums is closest, and does
$L_4 < \tfrac{16}{3} < R_4$ hold? Explain in one sentence why, using the fact
that $\sqrt{x}$ is increasing.
\end{exercise}

\begin{exercise}[An integral from the definition]
Following Example~\ref{ch20:ex-x2}, compute $\int_0^2 x\,\dd x$ directly from the
definition using a regular partition and right endpoints. You will need
$\sum_{k=1}^n k = \tfrac{n(n+1)}{2}$. Confirm that your answer equals the area of
the triangle you can see in the graph.
\end{exercise}

\begin{exercise}[Sum of cubes]
Using $\sum_{k=1}^n k^3 = \left(\tfrac{n(n+1)}{2}\right)^2$, evaluate
$\int_0^1 x^3\,\dd x$ from the definition with right endpoints. Show every step
of the limit.
\end{exercise}

\begin{exercise}[Using the properties]
Given $\int_1^4 f = 7$, $\int_1^4 g = -2$, and $\int_4^6 f = 5$, evaluate each of
the following or state that it cannot be determined:
(a) $\int_1^4 (3f - 2g)$;\quad
(b) $\int_1^6 f$;\quad
(c) $\int_6^4 f$;\quad
(d) $\int_1^4 f\cdot g$.
\end{exercise}

\begin{exercise}[Signed versus geometric area]
Consider $f(x) = x^2 - 1$ on $[0,2]$. (a) Sketch the graph and shade the regions
above and below the $x$-axis. (b) The signed area is $\int_0^2 (x^2-1)\,\dd x$;
using linearity, the known integral of $1$, and $\int_0^2 x^2\,\dd x = \tfrac{8}{3}$ (and $\int_0^1 x^2\,\dd x = \tfrac{1}{3}$), evaluate
it. (c) Now find the \emph{geometric} (total, unsigned) area by splitting at the
crossing point $x=1$ and adding absolute values. Explain why the two answers
differ.
\end{exercise}

\begin{exercise}[Average value]
Find the average value of $f(x) = x^3$ on $[0,2]$ (you may use
$\int_0^2 x^3\,\dd x = 4$). Then find a point $c \in [0,2]$ where $f(c)$ equals
that average value, confirming the Mean Value Theorem for integrals in this case.
\end{exercise}

\begin{exercise}[Accumulation reasoning]
Water flows into a tank at a rate $r(t) = 3 + 2t$ liters per minute, for $t$ from
$0$ to $5$ minutes. (a) Express the total volume that flows in as a definite
integral. (b) Without yet knowing antiderivatives, estimate it with a midpoint
sum using $5$ subintervals of width $1$. (c) The integrand is a straight line, so
the exact total is the area of a trapezoid -- compute it and compare to your
estimate.
\end{exercise}

\begin{exercise}[A conceptual trap]
A student claims: ``Since $\int_0^{2\pi}\sin x\,\dd x = 0$, the function $\sin x$
encloses no area over $[0,2\pi]$.'' Explain precisely what is right and what is
wrong with this statement, and write the correct integral expression for the
total geometric area enclosed between $\sin x$ and the $x$-axis on $[0,2\pi]$.
\end{exercise}

\begin{exercise}[Proof: monotonicity]
Prove the monotonicity property: if $f$ and $g$ are integrable on $[a,b]$ and
$f(x) \le g(x)$ for all $x \in [a,b]$, then $\int_a^b f \le \int_a^b g$.
\emph{Hint:} compare the two functions' Riemann sums on a common partition with
common sample points, term by term, then pass to the limit. Note where you use
$\Delta x_k > 0$.
\end{exercise}

\begin{exercise}[Proof: a bound from monotonicity]
Suppose $m \le f(x) \le M$ for all $x \in [a,b]$. Prove that
\[
  m\,(b-a) \;\le\; \int_a^b f(x)\,\dd x \;\le\; M\,(b-a).
\]
\emph{Hint:} apply the previous exercise twice, comparing $f$ to the constant
functions $m$ and $M$ and using $\int_a^b c\,\dd x = c(b-a)$. Then interpret the
result as a statement about the average value $\bar f$.
\end{exercise}

\begin{exercise}[Order of magnitude, no computation]
Without evaluating any integral, rank these three numbers from smallest to
largest and justify each comparison using monotonicity or a sketch:
$\int_0^1 x^2\,\dd x$,\quad $\int_0^1 x^3\,\dd x$,\quad $\int_0^1 \sqrt{x}\,\dd x$.
\emph{Hint:} on $(0,1)$, how do $x^3$, $x^2$, and $\sqrt{x}$ compare in size?
\end{exercise}
